% Options for packages loaded elsewhere
\PassOptionsToPackage{unicode}{hyperref}
\PassOptionsToPackage{hyphens}{url}
\PassOptionsToPackage{dvipsnames,svgnames,x11names}{xcolor}
\documentclass[
  english,
  10pt,
  a4paper,
]{article}
\usepackage{xcolor}
\usepackage[top=22mm,bottom=20mm,left=18mm,right=18mm]{geometry}
\usepackage{amsmath,amssymb}
\setcounter{secnumdepth}{-\maxdimen} % remove section numbering
\usepackage{iftex}
\ifPDFTeX
  \usepackage[T1]{fontenc}
  \usepackage[utf8]{inputenc}
  \usepackage{textcomp} % provide euro and other symbols
\else % if luatex or xetex
  \usepackage{unicode-math} % this also loads fontspec
  \defaultfontfeatures{Scale=MatchLowercase}
  \defaultfontfeatures[\rmfamily]{Ligatures=TeX,Scale=1}
\fi
\usepackage{lmodern}
\ifPDFTeX\else
  % xetex/luatex font selection
\fi
% Use upquote if available, for straight quotes in verbatim environments
\IfFileExists{upquote.sty}{\usepackage{upquote}}{}
\IfFileExists{microtype.sty}{% use microtype if available
  \usepackage[]{microtype}
  \UseMicrotypeSet[protrusion]{basicmath} % disable protrusion for tt fonts
}{}
\makeatletter
\@ifundefined{KOMAClassName}{% if non-KOMA class
  \IfFileExists{parskip.sty}{%
    \usepackage{parskip}
  }{% else
    \setlength{\parindent}{0pt}
    \setlength{\parskip}{6pt plus 2pt minus 1pt}}
}{% if KOMA class
  \KOMAoptions{parskip=half}}
\makeatother
\usepackage{longtable,booktabs,array}
\usepackage{caption}
\captionsetup[table]{skip=6pt}
\newcounter{none} % for unnumbered tables
\usepackage{calc} % for calculating minipage widths
% Correct order of tables after \paragraph or \subparagraph
\usepackage{etoolbox}
\makeatletter
\patchcmd\longtable{\par}{\if@noskipsec\mbox{}\fi\par}{}{}
\makeatother
\ifLuaTeX
\usepackage[bidi=basic,shorthands=off]{babel}
\else
\usepackage[bidi=default,shorthands=off]{babel}
\fi
\ifLuaTeX
  \usepackage{selnolig} % disable illegal ligatures
\fi
\setlength{\emergencystretch}{3em} % prevent overfull lines
\providecommand{\tightlist}{%
  \setlength{\itemsep}{0pt}\setlength{\parskip}{0pt}}
% ---- Çılgın Kurye document theme: same look as the CSE 396 project proposal ----
% Compile with XeLaTeX (also works on Overleaf). Fonts: TeX Gyre Termes (Times), DejaVu Sans Mono.
\usepackage{fontspec}
\setmainfont{texgyretermes}[Extension=.otf, UprightFont=*-regular, BoldFont=*-bold, ItalicFont=*-italic, BoldItalicFont=*-bolditalic]
\setsansfont{texgyreheros}[Extension=.otf, UprightFont=*-regular, BoldFont=*-bold, ItalicFont=*-italic, BoldItalicFont=*-bolditalic]
\setmonofont{DejaVuSansMono}[Extension=.ttf, UprightFont=*, BoldFont=*-Bold, ItalicFont=*-Oblique, BoldItalicFont=*-BoldOblique, Scale=0.80]
\setmathfont{texgyretermes-math.otf}

\usepackage{xcolor}
\definecolor{ckblue}{HTML}{1F3A5F}
\definecolor{cklight}{HTML}{E8EEF5}
\definecolor{ckgrey}{HTML}{F2F2F2}
\definecolor{ckline}{HTML}{555555}
\definecolor{shadecolor}{HTML}{F4F5F7}

% Headings as in the proposal: serif, bold; subsections bold italic. Numbers are part of the heading text.
\setcounter{secnumdepth}{0}
\usepackage{titlesec}
\titleformat{\section}{\rmfamily\Large\bfseries}{}{0pt}{}
\titleformat{\subsection}{\rmfamily\large\bfseries\itshape}{}{0pt}{}
\titleformat{\subsubsection}{\rmfamily\normalsize\bfseries}{}{0pt}{}
\titlespacing*{\section}{0pt}{12pt}{4pt}
\titlespacing*{\subsection}{0pt}{8pt}{3pt}

% Header and footer (the first page has only the footer, like the proposal)
\usepackage{fancyhdr,lastpage}
\pagestyle{fancy}
\fancyhf{}
\fancyhead[L]{\footnotesize\itshape CSE 396 · Group 8 · Çılgın Kurye}
\fancyhead[R]{\footnotesize\itshape\ckmodule}
\fancyfoot[C]{\footnotesize Page \thepage\ of \pageref*{LastPage}}
\renewcommand{\headrulewidth}{0.4pt}
\setlength{\headheight}{14pt}
\fancypagestyle{ckfirst}{\fancyhf{}\renewcommand{\headrulewidth}{0pt}\fancyfoot[C]{\footnotesize Page \thepage\ of \pageref*{LastPage}}}

% Inline code: breaks only after . : ( , / _ and spaces, never inside a name
\usepackage[obeyspaces,spaces]{url}
\AtBeginDocument{%
  \renewcommand{\UrlBreaks}{\do\.\do\:\do\(\do\,\do\/\do\_\do\|}%
  \renewcommand{\UrlBigBreaks}{}%
  \renewcommand{\UrlNoBreaks}{}%
  \urlstyle{tt}}

% Tables: boxed grid as in the proposal, smaller text
\usepackage{etoolbox}
\AtBeginEnvironment{longtable}{\footnotesize\renewcommand{\arraystretch}{1.2}}
\setlength{\LTpre}{6pt}
\setlength{\LTpost}{8pt}
\setlength{\arrayrulewidth}{0.4pt}

% Code blocks: long lines wrap; diagrams stay in one piece
\usepackage{framed}
\usepackage{fvextra}
\RecustomVerbatimEnvironment{verbatim}{Verbatim}{fontsize=\footnotesize, breaklines}
\makeatletter
\@ifundefined{Shaded}{}{\renewenvironment{Shaded}{\begin{snugshade}}{\end{snugshade}}}
\@ifundefined{Highlighting}{}{\RecustomVerbatimEnvironment{Highlighting}{Verbatim}{commandchars=\\\{\}, fontsize=\footnotesize, breaklines, breakanywhere}}
\makeatother

% Figures drawn with TikZ
\usepackage{tikz}
\usetikzlibrary{arrows.meta,positioning,fit,calc,backgrounds,shapes.geometric}
\tikzset{
  ck/.style={font=\footnotesize, >={Stealth[length=5pt]}, line width=0.6pt},
  box/.style={draw=ckline, fill=white, rounded corners=2pt, align=center, inner sep=4pt, minimum height=7mm},
  mod/.style={box, fill=ckgrey},
  this/.style={box, fill=cklight, draw=ckblue, line width=0.9pt},
  group/.style={draw=ckline, dashed, rounded corners=3pt, inner sep=6pt},
  lbl/.style={font=\scriptsize, align=center, inner sep=1.5pt},
  msg/.style={->, draw=ckline},
}
\usepackage{caption}
\captionsetup[figure]{font=small, labelfont=bf, skip=4pt}
\newcommand{\ckfigure}[2]{\begin{figure}[htbp]\centering #1\caption{#2}\end{figure}}

\usepackage{enumitem}
\setlist{itemsep=1pt, topsep=3pt}
\setlength{\parskip}{4pt}
\setlength{\parindent}{0pt}
\raggedbottom
% ---- end of theme ----
\newcommand{\ckmodule}{M8 LED Animation and Physical Board}
\usepackage{bookmark}
\IfFileExists{xurl.sty}{\usepackage{xurl}}{} % add URL line breaks if available
\urlstyle{same}
% fallback for those not using the hyperref driver hyperxmp:
\makeatletter
\@ifundefined{xmpquote}{\newcommand{\xmpquote}[1]{#1}}{}
\makeatother
\hypersetup{
  pdftitle={M8 LED Animation and Physical Board},
  pdflang={en},
  colorlinks=true,
  linkcolor={black},
  filecolor={Maroon},
  citecolor={Blue},
  urlcolor={black},
  pdfcreator={LaTeX}}

\title{M8 LED Animation and Physical Board --- Module Documentation}
\author{}
\date{}


% ---- colours and drawings ----
% ---- colours ----
\definecolor{srv}{HTML}{2E6FB7}\colorlet{srvL}{srv!15}
\definecolor{cli}{HTML}{2E9E5B}\colorlet{cliL}{cli!15}
\definecolor{adm}{HTML}{E08A1E}\colorlet{admL}{adm!18}
\definecolor{brd}{HTML}{7B4FB3}\colorlet{brdL}{brd!15}
\definecolor{tcpc}{HTML}{1F5FA8}
\definecolor{udpc}{HTML}{D9531E}
\definecolor{obs}{HTML}{D7263D}
\definecolor{tA}{HTML}{00A6D6}
\definecolor{tB}{HTML}{3BB273}
\definecolor{tC}{HTML}{F2C14E}
\definecolor{tD}{HTML}{9B5DE5}

% ---- styles ----
\tikzset{
  nb/.style={rounded corners=3pt, line width=0.9pt, align=center, inner sep=4pt, minimum height=9mm, font=\footnotesize},
  sv/.style={nb, draw=srv, fill=srvL},
  cl/.style={nb, draw=cli, fill=cliL},
  ad/.style={nb, draw=adm, fill=admL},
  bd/.style={nb, draw=brd, fill=brdL},
  hwb/.style={nb, draw=black!50, fill=black!5},
  tcpa/.style={draw=tcpc, line width=1.5pt},
  udpa/.style={draw=udpc, line width=1.5pt, dash pattern=on 5pt off 2.5pt},
  inna/.style={draw=black!55, line width=1pt},
  al/.style={font=\scriptsize, fill=white, inner sep=1.2pt, rounded corners=1pt},
  grp/.style={draw=black!35, dashed, rounded corners=5pt, inner sep=7pt},
  glbl/.style={font=\scriptsize\bfseries, text=black!65},
  hl/.style={line width=2.6pt},
  st/.style={nb, rounded corners=8pt, minimum width=24mm},
  co/.style={draw=black!60, line width=0.5pt},
  cot/.style={font=\scriptsize, align=left, text=black!80},
}

% ---- system map; argument = module to highlight (M1 ... M8, or none) ----
\newcommand{\systemmap}[1]{%
\tikzset{M1/.style={},M2/.style={},M3/.style={},M4/.style={},M5/.style={},M6/.style={},M7/.style={},M8/.style={},none/.style={}}%
\tikzset{#1/.style={hl}}%
\begin{tikzpicture}[ck]
\node[cl, M5, minimum width=27mm] (m5) at (0,3) {\textbf{M5}\\Player phones};
\node[ad, M6, minimum width=27mm] (m6) at (0,0.4) {\textbf{M6}\\Admin panel};
\node[sv, M2, minimum width=27mm, minimum height=38mm] (m2) at (4.6,1.7) {\textbf{M2}\\Network\\and sync};
\node[sv, M1, minimum width=27mm, minimum height=38mm] (m1) at (8.4,1.7) {\textbf{M1}\\Game server\\and rules};
\node[sv, M3, minimum width=25mm] (m3) at (12,2.8) {\textbf{M3}\\Physics};
\node[sv, M4, minimum width=25mm] (m4) at (12,0.6) {\textbf{M4}\\Traffic and AI};
\node[grp, fit=(m2)(m1)(m3)(m4)] (pc) {};
\node[glbl, anchor=south west] at (pc.north west) {MAIN COMPUTER};
\node[bd, M7, minimum width=27mm] (m7) at (4.6,-2.7) {\textbf{M7}\\ESP32 software};
\node[bd, M8, minimum width=27mm] (m8) at (8.4,-2.7) {\textbf{M8}\\LED engine};
\node[hwb, minimum width=25mm] (led) at (12,-2.7) {344 LEDs};
\node[hwb, minimum width=27mm] (sm) at (4.6,-4.4) {Sensor and\\4 motors};
\node[grp, fit=(m7)(m8)(led)(sm)] (bb) {};
\node[glbl, anchor=north east] at (bb.south east) {LED CITY BOARD};
\draw[tcpa, <->] (m5.east) -- node[al]{WebSocket} (m5.east -| m2.west);
\draw[tcpa, <->] (m6.east) -- node[al]{WebSocket} (m6.east -| m2.west);
\draw[tcpa, <->] ([xshift=-6mm]m7.north) -- node[al]{TCP} ([xshift=-6mm]m2.south);
\draw[udpa, <->] ([xshift=6mm]m7.north) -- node[al]{UDP} ([xshift=6mm]m2.south);
\draw[inna, <->] (m2) -- (m1);
\draw[inna, <->] (m1.east |- m3) -- (m3);
\draw[inna, <->] (m1.east |- m4) -- (m4);
\draw[inna, ->] (m7) -- node[al]{state} (m8);
\draw[inna, ->] (m8) -- (led);
\draw[inna, <->] (m7) -- (sm);
\draw[tcpa] (-1.3,-6.0) -- ++(0.9,0) node[right, font=\scriptsize]{TCP: every message arrives, in order};
\draw[udpa] (4.0,-6.0) -- ++(0.9,0) node[right, font=\scriptsize]{UDP: fast, the newest message wins};
\draw[inna] (9.3,-6.0) -- ++(0.9,0) node[right, font=\scriptsize]{inside one program or chip};
\end{tikzpicture}}

% ---- LED board seen from above (6 x 6 roads, 1 unit = 10 cm) ----
\newcommand{\boardbase}{%
\fill[black!88, rounded corners=3pt] (-0.5,-0.5) rectangle (5.5,5.5);
\foreach \i in {0,...,4}{\foreach \j in {0,...,4}{%
  \fill[black!60] (\i+0.13,\j+0.13) rectangle (\i+0.87,\j+0.87);
  \fill[black!22] (\i+0.24,\j+0.24) rectangle (\i+0.52,\j+0.52);
  \fill[black!30] (\i+0.56,\j+0.30) rectangle (\i+0.78,\j+0.62);
  \fill[black!38] (\i+0.28,\j+0.58) rectangle (\i+0.66,\j+0.78);}}
\foreach \i in {0,...,5}{%
  \draw[black, line width=4.5pt] (\i,-0.05) -- (\i,5.05);
  \draw[black, line width=4.5pt] (-0.05,\i) -- (5.05,\i);}
\foreach \i in {0,...,5}{\foreach \k in {0,...,30}{%
  \fill[black!55] (\i,\k/6) circle (0.025);
  \fill[black!55] (\k/6,\i) circle (0.025);}}
\foreach \x/\y/\c in {0/4/tA, 4/4/tB, 0/0/tC, 4/0/tD}{%
  \fill[\c!30!black!60] (\x+0.13,\y+0.13) rectangle (\x+0.87,\y+0.87);
  \fill[\c] (\x+0.3,\y+0.3) rectangle (\x+0.7,\y+0.7);
  \fill[white] (\x+0.4,\y+0.5) circle (0.05); \fill[white] (\x+0.6,\y+0.5) circle (0.05);}}
\newcommand{\nmark}[3]{\node[circle, fill=yellow!85!orange, draw=black, line width=0.4pt, inner sep=0.6pt, minimum size=3.6mm, font=\tiny\bfseries] at (#1,#2) {#3};}
\newcommand{\bike}[4]{%
\fill[#3!35] ($(#1,#2)+(#4:0.40)$) circle (0.05);
\fill[#3!65] ($(#1,#2)+(#4:0.21)$) circle (0.075);
\fill[#3, draw=white, line width=0.4pt] (#1,#2) circle (0.12);}
\newcommand{\dpoint}[3]{\draw[#3!50, line width=0.8pt] (#1,#2) circle (0.26); \draw[#3, line width=1.4pt] (#1,#2) circle (0.16); \fill[#3] (#1,#2) circle (0.07);}
\newcommand{\heavy}[2]{\draw[obs!60!black, line width=4pt] (#1) -- (#2); \draw[obs, line width=4pt, dash pattern=on 3pt off 5pt] (#1) -- (#2);}
\newcommand{\roadwork}[2]{\draw[obs, line width=4pt] (#1) -- (#2);}
\newcommand{\accident}[2]{\draw[obs, line width=4pt] (#1) -- (#2); \draw[white, line width=4pt, dash pattern=on 2.5pt off 3.5pt] (#1) -- (#2);}
\newcommand{\puddle}[2]{\draw[cyan!35, line width=4pt] (#1) -- (#2); \draw[white, line width=1.6pt, dash pattern=on 1pt off 2.5pt] (#1) -- (#2);}
\newcommand{\findme}[2]{\fill[white, opacity=0.25] (#1,#2) circle (0.36); \fill[white, opacity=0.5] (#1,#2) circle (0.24); \fill[white] (#1,#2) circle (0.13);}

% ---- phone game screen (7 x 3.4 units) ----
\newcommand{\phonemock}{%
\fill[black!85, rounded corners=9pt] (-0.3,-0.3) rectangle (7.3,3.7);
\fill[cyan!20] (0,1.7) rectangle (7,3.4);
\fill[cli!14] (0,0) rectangle (7,1.7);
\filldraw[fill=srv!30, draw=srv!70!black] (0.2,1.5) rectangle (1.0,3.0);
\filldraw[fill=adm!35, draw=adm!70!black] (1.1,1.6) rectangle (1.8,2.4);
\filldraw[fill=brd!30, draw=brd!70!black] (1.9,1.65) -- (2.6,1.65) -- (2.6,2.2) -- (2.25,2.5) -- (1.9,2.2) -- cycle;
\filldraw[fill=adm!35, draw=adm!70!black] (4.5,1.65) rectangle (5.1,2.3);
\filldraw[fill=srv!30, draw=srv!70!black] (5.2,1.55) rectangle (6.0,3.1);
\filldraw[fill=cli!35, draw=cli!70!black] (6.1,1.6) rectangle (6.85,2.6);
\fill[brown!60!black] (4.17,1.6) rectangle (4.23,1.8); \fill[cli!70!black] (4.2,1.95) circle (0.18);
\fill[black!55] (1.8,0) -- (5.2,0) -- (3.75,1.7) -- (3.25,1.7) -- cycle;
\draw[white, line width=0.8pt, dash pattern=on 4pt off 4pt] (3.5,0) -- (3.5,1.7);
\filldraw[fill=black!15, draw=black!70] (3.62,1.12) rectangle (3.88,1.32);
\fill[obs] (3.05,1.22) -- (3.15,1.48) -- (3.25,1.22) -- cycle;
\fill[black!85] (3.43,0.08) rectangle (3.57,0.3);
\fill[tA] (3.3,0.28) rectangle (3.7,0.6);
\fill[brown!65] (3.34,0.6) rectangle (3.66,0.82);
\fill[tA!60!black] (3.5,0.95) circle (0.13);
\node[fill=black!70, text=white, rounded corners=2pt, font=\scriptsize\bfseries, inner sep=2pt] at (3.5,3.15) {2:41};
\foreach \c/\s/\x in {tA/120/0.35,tB/95/0.95,tC/140/1.55,tD/80/2.15}
  \node[fill=\c, text=white, font=\tiny\bfseries, rounded corners=1pt, inner sep=1.5pt, minimum width=5mm] at (\x,3.15) {\s};
\fill[white, opacity=0.9] (6.45,2.85) circle (0.38); \draw[black!60] (6.45,2.85) circle (0.38);
\draw[-{Stealth[length=5pt]}, line width=1.8pt, obs] (6.45,2.85) ++(230:0.2) -- ++(50:0.5);
\fill[white, opacity=0.85] (0.55,0.5) circle (0.32); \fill[black!70] (0.43,0.5) -- (0.65,0.63) -- (0.65,0.37) -- cycle;
\fill[white, opacity=0.85] (1.4,0.5) circle (0.32); \fill[black!70] (1.52,0.5) -- (1.3,0.63) -- (1.3,0.37) -- cycle;
\node[fill=obs!80, text=white, rounded corners=3pt, font=\tiny\bfseries, minimum width=8mm, minimum height=6mm] at (5.6,0.5) {Brake};
\node[fill=cli, text=white, rounded corners=3pt, font=\tiny\bfseries, minimum width=8mm, minimum height=8mm] at (6.5,0.6) {Gas};
\node[fill=tC, text=black, rounded corners=3pt, font=\tiny\bfseries, minimum width=9mm, minimum height=5mm] at (6.15,1.45) {Deliver};
\node[fill=white, text=black, rounded corners=3pt, font=\tiny\bfseries, minimum width=9mm, minimum height=5mm] at (0.9,1.45) {Find me};}

% ---- side cut of the board (14 units wide) ----
\newcommand{\boardside}{%
\fill[brown!20] (-0.6,-0.35) rectangle (14.6,0);
\foreach \x/\h/\c in {1.0/0.7/black!18, 2.2/1.1/black!22, 4.6/1.9/black!25, 5.8/2.6/black!28, 7.4/2.2/black!25, 8.6/1.6/black!22, 10.9/1.0/black!18}
  \filldraw[fill=\c, draw=black!45] (\x,1.5) rectangle (\x+0.9,1.5+\h);
\filldraw[fill=tB!35, draw=tB!70!black] (12.4,1.5) rectangle (13.4,2.3);
\fill[black!70, rounded corners=1pt] (0.5,0) rectangle (1.1,0.9);
\fill[black!70, rounded corners=1pt] (12.9,0) rectangle (13.5,0.9);
\filldraw[fill=brown!45, draw=brown!70!black] (0,0.9) rectangle (14,1.3);
\fill[black!80] (0,1.3) rectangle (14,1.38);
\foreach \x in {0.3,0.9,...,13.8} \fill[black!35] (\x,1.42) circle (0.06);
\foreach \x/\c in {3.3/tA, 3.9/tA, 6.9/obs, 7.5/obs, 10.5/tC} \fill[\c] (\x,1.42) circle (0.07);
\fill[white, opacity=0.55] (0,1.38) rectangle (14,1.52);
\filldraw[fill=brd!40, draw=brd] (2.3,0.35) rectangle (4.3,0.85);
\node[font=\tiny] at (3.3,0.6) {ESP32};
\filldraw[fill=brd, draw=brd!60!black] (6.75,0.62) rectangle (7.25,0.88);
\filldraw[fill=adm!60, draw=adm!70!black] (1.8,0.6) circle (0.24);
\filldraw[fill=adm!60, draw=adm!70!black] (12.2,0.6) circle (0.24);
\filldraw[fill=black!15, draw=black!50] (9.3,0.35) rectangle (10.9,0.85);
\node[font=\tiny] at (10.1,0.6) {Power};}

\begin{document}

\thispagestyle{ckfirst}
\begin{center}
{\bfseries GEBZE TECHNICAL UNIVERSITY\par}
{Computer Engineering Department\par}
{CSE 396 Module Documentation\par}
\vspace{8pt}
{\LARGE\bfseries Çılgın Kurye\par}
\vspace{3pt}
{\large\itshape Module M8: LED Animation and Physical Board\par}
\vspace{6pt}
\textbf{Instructor:} Erkan Zergeroğlu\hspace{1.4em}\textbf{Company:} Board\&Byte\hspace{1.4em}\textbf{Group:} Group-8\hspace{1.4em}\textbf{Date:} 6 November 2026
\end{center}
\begin{center}\footnotesize\renewcommand{\arraystretch}{1.15}
\begin{tabular}{|p{0.27\linewidth}|p{0.15\linewidth}|p{0.27\linewidth}|p{0.15\linewidth}|}\hline
\textbf{Name} & \textbf{Student No.} & \textbf{Name} & \textbf{Student No.} \\ \hline
Ayşe Poyraz & 230104004085 & Hasan Hüseyin Özdemir & 230104004048 \\ \hline
Muhib Taha Boy & 230104004165 & Umutcan Ocak & 220104004903 \\ \hline
\end{tabular}
\end{center}
\vspace{2pt}



% ---- table and list helpers ----
\newcolumntype{L}[1]{>{\raggedright\arraybackslash}p{\dimexpr#1\linewidth-2\tabcolsep-\arrayrulewidth\relax}}
\newcolumntype{C}[1]{>{\centering\arraybackslash}p{\dimexpr#1\linewidth-2\tabcolsep-\arrayrulewidth\relax}}
\newcommand{\ckdone}{\item[$\boxtimes$]}
\newcommand{\ckopen}{\item[$\square$]}
\newcommand{\mdot}{$\bullet$}
\newcounter{msec}
\newcounter{chap}
% ---- end of helpers ----
\newcommand{\msection}[1]{\stepcounter{msec}\section{\arabic{msec}. #1}}

\msection{Purpose}

M8 builds the physical city board and writes the animation engine that turns the game state into
light. The board is 60 × 60 cm. It has 344 addressable LEDs on the roads and in the stores,
3D-printed buildings, and places for the motion sensor and the vibration motors. In the game the
board is the only map, so it must be clear to read from every side of the table.

\begin{longtable}[]{|L{0.18}|L{0.82}|}
\hline Runs on & Core 1 of the ESP32 for the software; the board itself for the hardware \\ \hline
Technology & Portable C++ animation engine that runs on the ESP32 and on a PC for tests; 344 addressable RGB LEDs; level shifter; 3D printing; plywood base \\ \hline
Requirements & R2 R5 R6 R7 R8 R11 \\ \hline
\end{longtable}

\begin{figure}[htbp]\centering
\systemmap{M8}
\caption{Where M8 sits in the system. M8 gets the state from M7 on the same chip.}
\end{figure}

\msection{Requirements covered}

\begin{longtable}[]{|L{0.07}|L{0.25}|L{0.68}|}
\hline \textbf{ID} & \textbf{Requirement} & \textbf{What M8 does for it} \\ \hline \endhead
R2 & Delivery cycle & Delivery points blink in the player's tone; a flash when the package is delivered \\ \hline
R5 & No map on the phone & The board is the map: all motorcycles; ``Find me'' white flash for 1 s \\ \hline
R6 & Obstacles & Red patterns by type on the whole road segment \\ \hline
R7 & Earthquake & A shock wave from the earthquake centre \\ \hline
R8 & Physical LED board & 60 × 60 cm, 6 roads in each direction, LEDs on all roads, printed buildings, four stores \\ \hline
R11 & Independent development & Demo mode without a server; animations checked as images on a PC \\ \hline
\end{longtable}

\msection{Inputs and outputs}

\begin{longtable}[]{|L{0.08}|L{0.20}|L{0.18}|L{0.54}|}
\hline \textbf{Dir.} & \textbf{Data} & \textbf{From / to} & \textbf{Content} \\ \hline \endhead
In & Board state & M7 & At most 30 per second: phase, players, delivery points, obstacles, earthquake \\ \hline
In & Config & M7 & Player colours, brightness 0.1--1, current limit \\ \hline
In & Link state & M7 & Connected, searching or lost \\ \hline
In & City plan & Shared file & Road grid and store places, turned into the LED map when the code is built \\ \hline
Out & Light & Board & 60 frames per second \\ \hline
Out & Frame report & M7 & Frame time and current estimate, shown to the admin \\ \hline
\end{longtable}

\msection{Main components}

\begin{longtable}[]{|L{0.20}|L{0.52}|L{0.28}|}
\hline \textbf{Component} & \textbf{What it does} & \textbf{Related modules} \\ \hline \endhead
Motion smoother & Moves the motorcycle dots smoothly between two states & M7 \\ \hline
LED map & The place of every LED; made from the city plan by a script & M1, city plan \\ \hline
Layers and effects & Five layers and the patterns for obstacles, earthquakes and game phases & --- \\ \hline
Post-processing and output & Colour correction, brightness, current limit; drives 5 LED lines at the same time & M7 \\ \hline
Demo mode & A built-in scenario that runs without a server & --- \\ \hline
\end{longtable}

\msection{How it works}

\textbf{The board.} The base is plywood. On it there are 6 roads in each direction, so the city has
5 × 5 blocks. LED strips with 60 LEDs per metre lie in the roads. A road segment between two
crossings is 10 cm long, so it has exactly 6 LED steps. With 336 LEDs on the roads and 2 LEDs in each
of the 4 stores, the board has 344 LEDs. The buildings are 3D-printed from the same models as the 3D
city on the phone, at a scale of 1 m in the game to 1 mm on the board. Thin printed covers over the
roads soften the light. The LEDs need 5 V. The software limits the current to 4 A, so the power
supply never gets too hot. A level shifter turns the 3.3 V signal of the ESP32 into a clean 5 V
signal.

\begin{figure}[htbp]\centering
\begin{tikzpicture}[x=1.05cm, y=1.05cm]
\boardbase
\bike{1}{3.5}{tA}{90}\bike{2.5}{4}{tA}{180}
\bike{3.5}{4}{tB}{0}\bike{5}{2.5}{tB}{90}
\bike{1.5}{1}{tC}{180}\bike{0}{1.5}{tC}{-90}
\bike{4}{2.5}{tD}{-90}\bike{4.5}{0}{tD}{180}
\dpoint{2.5}{2}{tA}\dpoint{5}{3.5}{tB}\dpoint{0.5}{3}{tC}\dpoint{3}{0.5}{tD}
\heavy{0,2}{1,2}
\roadwork{3,2}{3,3}
\accident{4,1}{5,1}
\puddle{2,4}{2,5}
\findme{3.5}{4}
\nmark{0.5}{5.25}{1}\nmark{0.62}{3.5}{2}\nmark{2.5}{2.45}{3}\nmark{0.5}{2.3}{4}\nmark{3.35}{2.5}{5}\nmark{4.5}{1.35}{6}\nmark{1.65}{4.5}{7}\nmark{3.5}{4.5}{8}
\node[anchor=west, font=\scriptsize, align=left] at (6.0,2.5) {\begin{tabular}{@{}c@{\ \ }l@{}}
\textbf{1} & Team store: two LEDs in team colour\\[2pt]
\textbf{2} & Motorcycle: a bright dot with a short tail\\[2pt]
\textbf{3} & Delivery point: blinks in the player's tone\\[2pt]
\textbf{4} & Heavy traffic: a slow dim red wave\\[2pt]
\textbf{5} & Road work: bright steady red\\[2pt]
\textbf{6} & Accident: red and white blinking\\[2pt]
\textbf{7} & Water puddle: a cool white shimmer\\[2pt]
\textbf{8} & ``Find me'': white flash for 1 second\end{tabular}};
\end{tikzpicture}
\caption{The planned LED board from above, during a game with 8 players.}
\end{figure}

\textbf{Colours.} There are four teams and each team has one colour: blue, green, yellow and purple.
The two players of a team have two tones of the team colour, a light one and a dark one. Red is
used only for obstacles, so red always means ``this road slows you down''. The team colours are far
from red, so nobody mixes them up. The pattern of the red tells the type of the obstacle, as the
figure below shows.

\begin{figure}[htbp]\centering
\begin{tikzpicture}[ck]
\foreach \r/\t in {0/Heavy traffic, 1/Road work, 2/Accident, 3/Water puddle}
  \node[font=\scriptsize\bfseries, anchor=east] at (2.3,-\r*0.75) {\t};
\foreach \k in {0,...,5}{%
  \pgfmathsetmacro{\a}{40+40*sin(\k*60)}
  \fill[obs!\a!black] (2.7+\k*0.32,0) circle (0.11);
  \fill[obs] (2.7+\k*0.32,-0.75) circle (0.11);
  \pgfmathtruncatemacro{\w}{mod(\k,2)}
  \ifnum\w=0 \fill[obs] (2.7+\k*0.32,-1.5) circle (0.11); \else \fill[white, draw=black!30] (2.7+\k*0.32,-1.5) circle (0.11); \fi
  \pgfmathsetmacro{\b}{30+20*sin(\k*97)}
  \fill[cyan!\b!white] (2.7+\k*0.32,-2.25) circle (0.11); \draw[black!20] (2.7+\k*0.32,-2.25) circle (0.11);}
\node[font=\scriptsize, anchor=west] at (4.75,0) {the light moves slowly along the road};
\node[font=\scriptsize, anchor=west] at (4.75,-0.75) {the whole road is red, it does not move};
\node[font=\scriptsize, anchor=west] at (4.75,-1.5) {red and white change twice per second};
\node[font=\scriptsize, anchor=west] at (4.75,-2.25) {soft white light that flickers};
\foreach \i/\c in {0/tA,1/tB,2/tC,3/tD}{%
  \fill[\c] (11.0+\i*0.75,-0.55) circle (0.17); \fill[\c!55] (11.0+\i*0.75,-1.05) circle (0.17);}
\node[font=\scriptsize\bfseries] at (12.12,0.0) {Team colours};
\node[font=\tiny] at (12.12,-1.5) {two tones in each team};
\end{tikzpicture}
\caption{Left: the six LEDs of one road segment for each obstacle type. Right: the team colours.}
\end{figure}

\textbf{Layers.} The engine draws 60 frames per second on core 1. Each frame is made from five
layers, from the bottom to the top: the roads, the obstacles, the delivery points, the motorcycles
and a top layer for short effects. A higher layer covers a lower one. The board state comes only 10
to 30 times per second, so the motion smoother moves each motorcycle dot between the last two
states. The dot has a short tail, so the direction of motion is easy to see. The top layer draws the
``Find me'' flash and the earthquake shock wave. After the layers, the engine corrects the colours
for the eye, sets the brightness from the config, checks the current and lowers all LEDs a little if
the current would be over the limit. Then it sends the frame to the LEDs on 5 lines at the same time,
so one frame takes only about 2 ms to send.

\begin{figure}[htbp]\centering
\begin{tikzpicture}[ck]
\foreach \i/\t/\c in {0/Roads/gray, 1/Obstacles/obs, 2/Delivery points/adm, 3/Motorcycles/tA, 4/Top effects/brd}{%
  \fill[\c!25, draw=\c, line width=0.9pt] (\i*0.35,\i*0.62) -- ++(4.2,0) -- ++(0.9,0.45) -- ++(-4.2,0) -- cycle;
  \node[font=\scriptsize, anchor=west] at (\i*0.35+5.3,\i*0.62+0.22) {\t};}
\draw[inna,->] (7.5,1.4) -- node[al]{mix} (8.8,1.4);
\node[nb, draw=srv, fill=srvL, minimum width=20mm] (pp) at (10.2,1.4) {Colour,\\brightness,\\current limit};
\draw[inna,->] (pp) -- (12.0,1.4);
\node[nb, draw=black!50, fill=black!8, minimum width=14mm] at (12.8,1.4) {5 LED\\lines};
\end{tikzpicture}
\caption{The five layers of one frame and the steps before the LEDs.}
\end{figure}

\textbf{Earthquake and link loss.} When an earthquake starts, a ring of light grows from the centre
of the earthquake over the whole board, like a wave on water, and slowly fades. The roads inside the
earthquake radius then turn red for road work. If M7 says that the link is lost, the roads go dark
and the store LEDs breathe slowly in amber, so everybody can see that the board is waiting.

\begin{figure}[htbp]\centering
\begin{tikzpicture}[x=0.62cm, y=0.62cm]
\foreach \s/\r in {0/0.6, 7.2/1.6, 14.4/2.8}{%
  \begin{scope}[shift={(\s,0)}]
  \boardbase
  \begin{scope}
  \clip (-0.5,-0.5) rectangle (5.5,5.5);
  \draw[white, line width=5pt, opacity=0.35] (3.4,2.6) circle (\r+0.25);
  \draw[yellow!70!orange, line width=3pt] (3.4,2.6) circle (\r);
  \end{scope}
  \fill[obs] (3.4,2.6) circle (0.12);
  \end{scope}}
\node[font=\scriptsize] at (2.5,-1.1) {start};
\node[font=\scriptsize] at (9.7,-1.1) {0.3 s later};
\node[font=\scriptsize] at (16.9,-1.1) {0.6 s later};
\end{tikzpicture}
\caption{The earthquake shock wave grows from the centre of the hit.}
\end{figure}

\msection{Testing and demonstration}

The engine is plain C++, so it also runs on a PC. There, every animation is drawn as an image and
compared with a saved reference image. On the real board, built-in test patterns check the LED order
and place, the 8 player colours (three people must name each colour correctly), the current limit,
the voltage at the farthest LED, and 30 minutes of heat. In the demonstration the board runs its demo
scenario without a server, with all obstacle types and an earthquake.

\end{document}
